\section{Introduction}
\label{sec:intro}

4D scene generation seeks to reconstruct comprehensive spatiotemporal representations capturing both explicit 3D geometry and complex temporal dynamics. Generating such dynamically rich 4D scenes from \textit{a single static image} remains a fundamental challenge~\cite{miao2025advances}, as the model must infer plausible future motion from a static observation while preserving structural coherence across time and viewpoints. Achieving this capability would largely benefit applications such as virtual reality, augmented reality~\cite{guo2023vid2avatar,miao2024pla4d}, and immersive content creation.

Although recent video generation models can produce realistic dynamic content, they generally lack an explicit understanding of 3D structure, leading to view inconsistencies and failing to capture physically plausible motion. Existing single-image-to-4D methods therefore typically adopt one of two decoupled paradigms. \textit{Generate-then-reconstruct} methods~\cite{cat4d,4real,dimensionX,free4d} use powerful video models to synthesize multi-view videos before reconstructing a 4D representation. This paradigm benefits from the high-fidelity and rich dynamics offered by existing generation models. 
However, video models struggle to maintain strict geometric consistency across the generated views, leading to significant artifacts and structural collapse during the subsequent 3D reconstruction phase. Consequently, the alternative paradigm, \textit{reconstruct-then-generate}~\cite{animate124,animate3d,gs-dit,jin2025optimizing, gen3c}, has emerged as another promising direction. By first establishing a static 3D structure, this approach provides a robust geometric foundation for subsequent motion generation. 
Still, by decoupling static geometry from motion generation, they discard the rich dynamic potential latent in the source image. As a result, they are restricted to modeling physically plausible and externally constrained motions (e.g., swinging) and struggle to generate large-scale and self-initiated movements that originate from the scene or objects themselves.

To address these challenges, we revisit the \textit{reconstruct-then-generate} paradigm and make its interface between structure and dynamics substantially tighter, as illustrated in Figure~\ref{fig:teaser}. Rather than explicitly refining the base geometry during generation, we condition future dynamic generation on an initial geometric prior inferred from the input image, and model the scene as dense 4D point trajectories, which is directly predicted from the single input image with an integrated model. To facilitate training of the 3D motion generation model, we first built TrajScene-60K, a large-scale dataset of 60,000 samples with 4D point trajectories, tackling data scarcity of large-scale, high-quality 4D scene data with complex dynamics. We then introduce the \textit{4D Scene Trajectory Generator} (4D-STraG) that takes an input image and predicts future time-varying point trajectories relative to the initial frame. Unlike prior works that treat these as separate steps, 4D-STraG operationalizes them within a unified denoising process, producing coherent 4D point trajectories whose motion is intrinsically consistent with the evolving 3D structure. To fully leverage priors from the input image, we further incorporate a depth-guided motion normalization method to enhance geometric awareness, and a Motion Perception Module (MPM) to leverage plausible motion priors. Finally, the generated 4D point cloud trajectories are rendered into high-fidelity dynamic videos from arbitrary novel viewpoints using our \textit{4D View Synthesis Module} (4D-ViSM), completing the full pipeline from a single image to a coherent 4D scene. 

Experimental results of both quantitative and qualitative comparisons demonstrate that our method consistently outperforms existing baselines, producing 4D scenes with more pronounced dynamics, stronger 3D motion consistency, and superior performance. Detailed ablation studies further confirm the effectiveness of our key modules, highlighting how each contributes to coherent and physically plausible 4D motion. Our contributions can be summarized as:

\begin{itemize}
    \item We construct TrajScene-60K, a large-scale 4D scene dataset featuring dense 4D point cloud trajectories, videos, and text to advance research in this field.
    \item We propose a geometry-conditioned dense-trajectory formulation for single-image 4D synthesis. It predicts scene trajectories from an initial geometric prior, avoiding the reliance on loosely decoupled stages.
    \item We develop 4D-STraG, a diffusion-based trajectory generator with depth-guided motion normalization and a Motion Perception Module for structural stability and motion plausibility, alongside 4D-ViSM for high-quality novel-view dynamic rendering from the generated 4D representation.
    \item Extensive experiments show that \ours delivers strong perceptual quality, geometry-aware consistency, and efficient 4D generation.
\end{itemize}